\documentclass{article}
\usepackage[utf8]{inputenc}
\usepackage[T1]{fontenc}
\usepackage{url}
\usepackage{graphicx}
\usepackage{geometry}
\usepackage{listings}
\usepackage{hyperref}
\usepackage{multicol}
\usepackage{array}
\usepackage{booktabs}
\usepackage{xcolor}
\usepackage{fancyhdr}
\usepackage{titlesec}
\usepackage{longtable}
\usepackage[most]{tcolorbox}

\lstset{numbers=none, basicstyle=\ttfamily}
\renewcommand{\lstlistingname}

\geometry{a4paper, left=20mm, right=20mm, top=20mm, bottom=20mm}


% ==========================================
% Colors
% ==========================================

\definecolor{codebg}{RGB}{245,245,245}
\definecolor{codegreen}{RGB}{0,120,0}
\definecolor{codeblue}{RGB}{0,0,180}
\definecolor{codered}{RGB}{180,0,0}

% ==========================================
% Listings Configuration
% ==========================================

\lstset{
backgroundcolor=\color{codebg},
basicstyle=\ttfamily\small,
breaklines=true,
frame=single,
columns=fullflexible,
keywordstyle=\color{codeblue}\bfseries,
commentstyle=\color{codegreen},
stringstyle=\color{codered},
showstringspaces=false
}

% ==========================================
% Custom Boxes
% ==========================================

\newtcolorbox{tipbox}{
colback=green!5,
colframe=green!60!black,
title=Tip
}

\newtcolorbox{warningbox}{
colback=red!5,
colframe=red!60!black,
title=Important
}

\newtcolorbox{notebox}{
colback=blue!5,
colframe=blue!60!black,
title=Note
}

% ==========================================
% Document
% ==========================================

\begin{document}

\begin{titlepage}
\centering

\author{Eyad Islam El-Taher}
{\Huge\bfseries Open URL Redirect Notes \par}

\vspace{0.5cm} 

{\Large Author: Eyad Islam El-Taher \par}

\vspace{0.5cm}

{\large \today \par}

\vfill

\end{titlepage}





% ============================================================
% Introduction
% ============================================================

\section{Introduction}

Open URL Redirect (also known as Open Redirect or Unvalidated Redirect) is a web application vulnerability that occurs when an application accepts user-controlled input and uses it as the destination of a redirect without proper validation.

This vulnerability allows attackers to craft URLs that appear to belong to a trusted domain while redirecting victims to attacker-controlled websites. Open Redirect vulnerabilities are commonly abused in phishing attacks, OAuth attacks, token theft scenarios, and as part of larger attack chains.

\begin{notebox}
Although Open Redirect is often classified as a low-severity vulnerability, its impact can become critical when chained with authentication mechanisms, OAuth flows, or social engineering attacks.
\end{notebox}

% ============================================================
% HTTP Redirection Status Codes
% ============================================================

\section{HTTP Redirection Status Codes}

HTTP redirection responses belong to the 3xx class of status codes.

\begin{longtable}{p{2cm}p{11cm}}
\toprule
\textbf{Code} & \textbf{Description} \\
\midrule
300 & Multiple Choices \\
301 & Moved Permanently \\
302 & Found (Temporary Redirect) \\
303 & See Other \\
304 & Not Modified \\
305 & Use Proxy \\
307 & Temporary Redirect \\
308 & Permanent Redirect \\
\bottomrule
\end{longtable}

\begin{notebox}
Open Redirect vulnerabilities are most commonly observed in HTTP responses returning status codes 301, 302, 307, and 308 with a user-controlled \texttt{Location} header.
\end{notebox}

% ============================================================
% Redirect Methods
% ============================================================

\section{Redirect Methods}

Applications may implement redirection using several methods.

\subsection{HTTP Header Redirect}

\begin{lstlisting}[language=PHP]
header("Location: https://example.com");
exit();
\end{lstlisting}

\subsection{Path-Based Redirect}

\begin{lstlisting}
https://example.com/redirect/http://evil.com
https://example.com/redirect/../http://evil.com
\end{lstlisting}

\subsection{JavaScript-Based Redirect}

\begin{lstlisting}
window.location = redirectUrl;
\end{lstlisting}

\subsection{Meta Refresh Redirect}

\begin{lstlisting}[language=HTML] <meta http-equiv="refresh"
content="0;url=https://example.com">
\end{lstlisting}

% ============================================================
% Why Open Redirect Happens
% ============================================================

\section{Why Open Redirect Happens}

Open Redirect vulnerabilities occur when developers trust user-controlled input and use it directly in redirection logic.

Common causes include:

\begin{itemize}
\item Lack of URL validation.
\item Missing allowlists.
\item Weak regular expressions.
\item OAuth misconfigurations.
\item Trusting GET or POST parameters.
\item Client-side JavaScript redirects.
\end{itemize}

% ============================================================
% Vulnerable Code Examples
% ============================================================

\section{Vulnerable Code Examples}

\subsection{PHP Native}

\begin{lstlisting}

<?php

$url = $_GET['url'];

header("Location: " . $url);
exit();

?>

\end{lstlisting}

\subsection{Laravel}

\begin{lstlisting}
public function redirect(Request $request)
{
return redirect($request->url);
}
\end{lstlisting}

\subsection{JavaScript}

\begin{lstlisting}
const redirectUrl =
new URLSearchParams(window.location.search)
.get("url");

window.location.href = redirectUrl;
\end{lstlisting}

\begin{warningbox}
Never use user-controlled input directly in redirect functions.
\end{warningbox}

% ============================================================
% Where To Find Open Redirects
% ============================================================

\section{Where to Find Open Redirects}

Common locations include:

\begin{itemize}
\item Login pages
\item Logout pages
\item OAuth implementations
\item SSO systems
\item Password reset workflows
\item Email verification pages
\item Checkout systems
\item Payment gateways
\item Affiliate programs
\item Tracking links
\item Redirect endpoints
\item Mobile deep links
\end{itemize}

% ============================================================
% Common Parameter Names
% ============================================================

\section{Common Parameter Names}

\begin{lstlisting}
url=
redirect=
redirect_url=
redirect_uri=
redir=
next=
continue=
return=
return_to=
returnTo=
destination=
dest=
target=
go=
rurl=
checkout_url=
view=
\end{lstlisting}

% ============================================================
% Payloads
% ============================================================

\section{Payloads}

Common testing payloads:

\begin{lstlisting}
https://evil.com
http://evil.com
//evil.com
////evil.com
https:evil.com
//evil.com
/%2f%2fevil.com
https://attacker.com
\end{lstlisting}

% ============================================================
% DOM-Based Open Redirect Example
% ============================================================

\section{DOM-Based Open Redirect Example}

Consider the following vulnerable JavaScript code:

\begin{lstlisting}
<a href='#' onclick='returnURL =
/url=(https?://.+)/.exec(location);

if(returnUrl)
location.href = returnUrl[1];
else
location.href = "/"'>
Back to Blog </a>
\end{lstlisting}

The application extracts the value of the \texttt{url} parameter from the current URL and directly assigns it to \texttt{location.href}.

Example:

\begin{lstlisting}
https://victim.com/post?postId=4&
url=https://evil.com
\end{lstlisting}

When the victim clicks \texttt{Back to Blog}, the browser redirects to the attacker-controlled website.

\begin{notebox}
This vulnerability exists entirely on the client side and is therefore classified as a DOM-Based Open Redirect.
\end{notebox}

% ============================================================
% Methodology
% ============================================================

\section{Methodology}

\begin{enumerate}
\item Identify parameters that influence navigation.
\item Replace values with an external domain.
\item Observe redirection behavior.
\item Analyze HTTP response headers.
\item Test protocol-relative URLs.
\item Test encoded payloads.
\item Test JavaScript-based redirects.
\item Test OAuth redirect URIs.
\item Attempt filter bypass techniques.
\end{enumerate}

% ============================================================
% Filter Bypass Techniques
% ============================================================

\section{Filter Bypass Techniques}

\subsection{Whitelisted Domain Bypass}

\begin{lstlisting}
[www.whitelisted.com.evil.com](http://www.whitelisted.com.evil.com)
\end{lstlisting}

\subsection{Protocol Relative URLs}

\begin{lstlisting}
//evil.com
////evil.com
\end{lstlisting}

\subsection{Bypass Blacklisted Slash}

\begin{lstlisting}
https:evil.com
\end{lstlisting}

\subsection{Escaped Slash Bypass}

\begin{lstlisting}
//evil.com
//evil.com
\end{lstlisting}

\subsection{Null Byte Injection}

\begin{lstlisting}
[https://evil.com%00](https://evil.com%00)
\end{lstlisting}

\subsection{HTTP Parameter Pollution}

\begin{lstlisting}
?next=trusted.com&next=evil.com
\end{lstlisting}

\subsection{Username@Host Trick}

\begin{lstlisting}
https://trusted.com@evil.com
\end{lstlisting}

\subsection{Unicode Normalization}

\begin{lstlisting}
https://evil.c[U+2100].example.com
\end{lstlisting}

% ============================================================
% Impact
% ============================================================

\section{Impact}

Possible impacts include:

\begin{itemize}
\item Phishing attacks
\item Credential theft
\item OAuth token theft
\item Access control bypass
\item Session theft
\item Malware distribution
\item Security filter bypass
\item Brand reputation damage
\end{itemize}

% ============================================================
% Remediation
% ============================================================

\section{Remediation}

Recommended mitigations:

\begin{itemize}
\item Use allowlists for redirect destinations.
\item Restrict redirects to internal paths.
\item Validate hosts and schemes.
\item Reject protocol-relative URLs.
\item Validate OAuth redirect URIs exactly.
\item Avoid user-controlled redirects whenever possible.
\item Use indirect reference identifiers.
\end{itemize}

% ============================================================
% References
% ============================================================

\end{document}